\chapter{Przestrzenie topologiczne.}

\section{Topologia.}

\begin{definition}
	Niech $X$ będzie nie pustym zbiorem. Mówimy, że podzbiór $\mathcal{T}\subseteq X $ jest \textbf{topologią na $X$} jeśli:
    \begin{enumerate}[label=\arabic*)]
        \item Zbiór $X$ oraz zbiór pusty $\emptyset$, należą do $\mathcal{T}$.
        \item Zbiór $\mathcal{T}$ jest zamknięty na branie sumy (skończona lub nieskończone).
        \item Zbiór $\mathcal{T}$, jest zamknięty na branie przekroju.    
       \end{enumerate}

  Parę $\left( X, \mathcal{T} \right)  $ nazywamy \textbf{przestrzenią topologiczną.}   
\end{definition} 

\begin{definition}
    Niech $X$ będzie niepustym zbiorem oraz niech $\mathcal{T}$ będzie rodziną wszystkich podzbiorów $X$. Wtedy $\mathcal{T}$ jest topologią, którą nazywamy \textbf{topologią dyskretną}. 

Przestrzeń topologiczna $\left( X, \mathcal{T} \right)$ nazywa się wtedy \textbf{przestrzenią dyskretna}.         
\end{definition} 
\begin{proof}
    Dowód, że tak zdefiniowane $\mathcal{T}$ jest topologią. 
\noindent$\bullet$ $\emptyset \subseteq X $ oraz $X \subseteq X $, więc $\emptyset, X \in \mathcal{T}$.
\noindent$\bullet$ Rozważmy dowolną rodzinę zbiorów $\{X_{i} \}_{i \in I} \in \mathcal{T}$. Ponieważ dla dowolnego $i$, $X_{i} \subseteq X$, więc $\bigcup_{i \in I} X_{i}  \subseteq X $, a co za tym idzie $\bigcup_{i \in I} X_{i} \in \mathcal{T}$.
\noindent$\bullet$ Rozważmy dwa dowolne zbiory $A, B \in \mathcal{T}$. Wtedy $A\cap B \subseteq X $, a co za tym idzie $A\cap B \in \mathcal{T}.$  
\end{proof}

\begin{definition}
    Niech $X$ będzie niepustym zbiorem oraz niech $\mathcal{T} = \{\emptyset, X\}$. Wtedy $\mathcal{T}$ nazywamy \textbf{topologią niedyskretną} oraz parę $\left( X, \mathcal{T} \right)  $ \textbf{przestrzenią niedyskretną}.      
\end{definition}  
\begin{proof}
    Oczywiście $\emptyset$ oraz $X$ należą do $\mathcal{T}$. Dowolna suma jak i dowolny przekrój może dać w wyniku tylko $\emptyset$ lub $X$, które z definicji są w $\mathcal{T}$. Zatem $\mathcal{T}$ jest topologią.        
\end{proof} 

\begin{theorem}
    Jeśli $\left( X, \mathcal{T} \right)  $ jest przestrzenią topologiczną, taką że dla dowolnego $x \in X$, singleton $\{x\} \in \mathcal{T} $, to $\mathcal{T}$ jest topologią dyskretną. 
\end{theorem} 
\begin{proof}
    Jeśli $S$ jest dowolnym podzbiorem $X$, to możemy zapisać $S$ jako sumę
   \[
       S = \bigcup_{x \in S} \{x\}    
   .\]
    Zatem w $\mathcal{T}$ jest zawarty dowolny podzbiór $X$, czyli $\mathcal{T}$ jest topologią dyskretną. 
\end{proof} 
\subsubsection*{Ćwiczenia.}
\begin{problem}[5]
    Udowodnij, że każda z poniższych rodzin podzbiorów $\mathbb{R}$ jest topologią.
\end{problem} 
\subsection{Zbiory otwarte, domknięte i clopen.}

\[
    \begin{tikzcd} \mathbb{F}_X \arrow[r, "\overline{f}"] \arrow[d, "\pi"'] & H \\ G = \mathbb{F}_X/N \arrow[ru, "\varphi"'] \end{tikzcd}
\]
